\documentstyle[%


righttag%


,11pt%


,amssymb%


,russian%               remove in Eng.


,izvru%                 Set izven% in Eng.


]{amsart}





%       Math definitions





\newcommand{\tts}[1]{}





%\hoffset=-15mm%                  For Eng.


\setcounter{page}{74}


\god{1999}


\nomer{12~(451)} %                        Remove in Eng.


%\nomer{Vol.\,43, No.\,12}%               Set in Eng.


%\str{\pageref{begin}--\pageref{end}}%  Set in Eng.


\udk{519.854}





\author{\it А.С.\,Стрекаловский, А.А.\,Кузнецова}


% NB:   ^^ remove in Eng.


\title{О сходимости алгоритма глобального поиска  в задаче


выпуклой максимизации на допустимом множестве}


\thanks{Работа выполнена при


финансовой поддержке Российского фонда фундаментальных


исследований (грант \No\,98-01-00043).}





\date{}


%\pagestyle{myheadings}%         Set in Eng.


\pagestyle{plain} %              Remove in Eng.


\begin{document}


%\thispagestyle{plain}%          Set in Eng.


\maketitle


\label{begin}








\section{Введение}





В последние годы невыпуклые задачи оптимизации обращают на себя все


большее внимание, в частности, ввиду широкого поля приложений 


\cite{l1}--\cite{l5}. Примером является задача выпуклой максимизации 


\cite{l3}--\cite{l7}


\begin{equation}


f(x)\rightarrow \max,\quad x\in D,


\label{eq:m1} 


\end{equation}


где $f:R^{n}\rightarrow R\cup \{+\infty \}$ и $D\subset R^{n}$ выпуклы.


Несмотря на внешнюю простоту, эта задача оказывается сложной


как для исследований, так и, главное, для численного решения, 


поскольку в ней может существовать большое количество


локальных максимумов, не являющихся глобальными \cite{l1}--\cite{l7}.


Поэтому для подобных задач считается естественным применение 


популярных методов ветвей и границ, отсечений, погружений и т.\,п.


В отличие от этих подходов в работах \cite{l6}--\cite{l10}


методы глобального поиска основаны на теории необходимых 


и достаточных условий глобальной оптимальности. 





Естественным выглядит вопрос о глобальной сходимости 


таких алгоритмов. 


Теоремы сходимости дают, по-видимому, определенную теоретическую поддержку 


численным экспериментам (\cite{l1}, с.\,415--420; \cite{l2}, с.\,40--42).


Поэтому основной целью данной статьи является доказательство теоремы


сходимости алгоритма глобального поиска, предлагаемого ниже, для


задачи (1). Для этого используется вводящееся здесь понятие разрешающего


набора.


Ранее \cite{l7} подобная теорема была доказана только для квадратичной 


задачи. К тому же некоторые требования из \cite{l7} удалось


значительно ослабить. Для задачи (1) доказана также сходимость


метода локального подъема.





\section{Метод локального подъема}





В дальнейшем будем считать, что в задаче (\ref{eq:m1}) функция


$f:R^{n}\rightarrow R$ является 


выпуклой и дифференцируемой в 


некоторой (открытой) области $Q$, содержащей допустимое множество


и множество Лебега функции $f(\cdot)$ вида


$L(f)=\{x\in R \mid f(x)\leq \sup(f,D) \}$. При этом предполагается,


что функция $f(\cdot)$ ограничена сверху на $D$


\begin{equation}


\sup (f,D)\triangleq f_{*}<+\infty.


\label{eq:m2}


\end{equation}


Как известно (\cite{l2}, теорема~1, с.\,179), условие 


\begin{equation}


\langle \nabla f(z), x-z \rangle \leq 0 \ \; \forall x \in D 


\label{eq:m3}


\end{equation}


локального максимума


в задаче (\ref{eq:m1}) не является достаточным для глобальности


максимума в


допустимой точке $z\in D$.


Эта точка является стационарной или критической (т.\,е. удовлетворяющей


условию (\ref{eq:m3})), если она является решением линеаризованной 


в точке $z$ задачи


$$


\langle \nabla f(z), x \rangle \uparrow \max, \ \; x \in D.


$$ 





Приведем простейшую вычислительную процедуру, основанную на решении


линеаризованной задачи. Пусть дана числовая последовательность 


$\{\delta_{s}\}$, $\delta_{s}>0$, $s=0,1,2,\dotsc$, и 


допустимая точка $p^{0}\in D$. Последовательность $\{p^{s}\}$ будем строить


по правилу


\begin{equation}


p^{s+1}\in D : \; \langle \nabla f(p^{s}), p^{s+1} \rangle + \delta_{s} \geq 


\sup_{x \in D} \langle \nabla f(p^{s}), x \rangle . 


\label{eq:m4}


\end{equation} 


Оказывается, такая простая вычислительная процедура при сделанных


предположениях в задаче (\ref{eq:m1}) генерирует последовательность,


сходящуюся к стационарной точке в следующем смысле.





\begin{prop}


Пусть функция $f$ выпукла и непрерывно дифференцируема на 


открытой области $Q$ и ограничена сверху на допустимом множестве $D$,


и пусть


$\sum\limits_{s=1}^{\infty}\delta_{s}<+\infty$. Тогда


последовательность $\{p^{s}\}$, генерируемая по правилу $(\ref{eq:m4})$,


удовлетворяет условию


\begin{equation}


\lim_{s\rightarrow +\infty}


[\sup_{x \in D} \langle \nabla f(p^{s}), x - p^{s}\rangle]=0 . 


\label{eq:m5}


\end{equation}


 


В случае ограниченности множества $D$ любая предельная точка $z$


последовательности $\{p^s\}$ удовлетворяет условию оптимальности


$(\ref{eq:m3})$. 


\end{prop}





{\bf Доказательство.}


Из (\ref{eq:m4}) с учетом выпуклости функции $f(\cdot)$ следует


\begin{equation}


0\leq\sup_{x \in D} \langle \nabla f(p^{s}), x - p^{s}\rangle\leq 


\langle \nabla f(p^{s}), p^{s+1} - p^{s}\rangle + \delta_{s}\leq


f(p^{s+1})-f(p^{s})+\delta_{s}. 


\label{eq:m6}


\end{equation} 


Отсюда


$$


f(p^{s})\leq f(p^{s+1})+ \delta_{s}.


$$


Последнее неравенство вместе с (\ref{eq:m2}) и условием


$\sum\limits_{s=1}^{\infty}\delta_{s} < +\infty $


гарантирует ({\cite{l1}, лемма~2.3.2}), что числовая последовательность


$\{f(p^{s})\}$ сходится. Поэтому из (\ref{eq:m6}) следует (\ref{eq:m5}).


А тогда последнее утверждение предложения становится очевидным.





\begin{note}


Как известно (напр., \cite{l1}, с.\,291--298),


процедура из


(\ref{eq:m4}) входит в стандартные градиентные методы выпуклой оптимизации


для нахождения вспомогательной точки $\ov{x}^{k}$, с помощью которой затем,


скажем, составлением выпуклой комбинации 


$x(\alpha)=\alpha \ov{x}^{k}+(1-\alpha)x^{k}$


поиск следующей точки $x^{k+1}$ сводится к одномерной минимизации функции


$\varphi (\alpha)\triangleq f(x(\alpha))$ по параметру $\alpha$.


В методе (\ref{eq:m4}) сложный одномерный поиск исключается,


результат оказывается тот же --- стационарная точка. В выпуклой задаче


стационарная точка уже будет глобальным решением, в то время как 


в задаче (\ref{eq:m1}) это не так.


\end{note}





\begin{note}


При реальном счете может быть использовано несколько 


критериев останова метода (\ref{eq:m4}) от очень жесткого


$$


\| p^{s+1}-p^{s} \| \leq \frac{\varepsilon}{2\| \nabla f(p^{s}) \|}


$$


до более слабых


$$


f(p^{s+1})-f(p^{s})\leq \varepsilon/2,\quad


\langle \nabla f(p^{s}), p^{s+1} - p^{s}\rangle \leq 


\varepsilon/2,


$$


где $\varepsilon$ --- заданная точность. При этом из (\ref{eq:m6}) легко


получить, что при $\delta_{s} \leq \varepsilon/2$


точка $p^{s}$ оказывается $\varepsilon$-критической.


\end{note}





\section{Стратегия глобального поиска}





Для задачи (\ref{eq:m1}) справедливы следующие


необходимые и достаточные условия глобальности решения.





\begin{thm}[{\series{m}\shape{n}\selectfont\cite{l6}, \cite{l7}}]


\label{Theorem1}


Пусть $D$ выпукло, $z \in D$ и, кроме того,


\begin{equation}


-\infty \leq \inf (f,D)< f(z)< +\infty.


\label{eq:m8}


\end{equation}


Тогда, для того чтобы $z$ было глобальным решением задачи $(\ref{eq:m1})$,


необходимо и достаточно, чтобы


\begin{equation}


\begin{array}{c}


\forall y \in D: \; f(y)=f(z),\\[2mm]


\langle \nabla f(y), x-y \rangle \leq 0 \ \; \forall x \in D.


\end{array}


\label{eq:m9}


\end{equation}


\end{thm} 





\begin{note}


Нетрудно видеть, что условие (\ref{eq:m9})


связано с классической теорией экстремума. Полагая $y=z$ в (\ref{eq:m9}),


получаем (\ref{eq:m3}). Отметим, что использование в условии (\ref{eq:m9})


только допустимых $y$ было предложено в \cite{l8}. С другой стороны,


доказательство теоремы~\ref{Theorem1} \cite{l6}, \cite{l7}


демонстрирует так называемое


алгоритмическое свойство условия оптимальности (\ref{eq:m9}),


т.\,е. если условие оптимальности нарушено, то можно построить


допустимую точку лучшую, чем исследуемая. Нетрудно видеть, что все


классические условия оптимальности (в частности (\ref{eq:m3})) обладают


алгоритмическим свойством. Используя это свойство, 


можно предложить алгоритм глобального поиска в задаче (\ref{eq:m1}). 


\end{note}





С помощью функции


$$


\varphi (z)=\sup_{x,y}\{\langle \nabla f(y),x-y\rangle \mid


x,y \in D,\ f(y)=f(z)\}


$$


условие (\ref{eq:m9}) можно преобразовать \cite{l7} к виду


$\varphi (z) \leq 0$.


Иначе говоря, для проверки этого условия необходимо уметь решать


задачу


\begin{equation}


\langle \nabla f(y),x-y\rangle \uparrow \sup_{x,y},\ \;


x,y\in D,\ \; f(y)=f(z).


\label{eq:m11}


\end{equation}





Поскольку задача совместной по $x$ и $y$ максимизации представляется


достаточно трудной, предлагается ее разбить на три более простые


задачи. Первая --- задача максимизации по $x$


\begin{equation}


\langle \nabla f(y),x\rangle \uparrow \max_{x}, \ \; x\in D,


\label{eq:m12}


\end{equation}


где точка $y$ такая, что $f(y)=f(z)$,


будет называться линеаризованной задачей, а вторая --- 


максимизация по $v$ при фиксированном $u\in D$


\begin{equation}


\langle \nabla f(v),u-v\rangle \uparrow \max_{v}, \ \; f(v)=f(z),


\label{eq:m13}


\end{equation}


--- задачей уровня.


Поэтому совершенно естественными выглядят следующие предположения.





\begin{itemize}


\item[(HL):] для любого $\varepsilon >0$ и для любого


$y \in D: f(y)\geq \inf(f,D)$


можно приближенно решить линеаризованную задачу (\ref{eq:m12}),


т.\,е. отыскать элемент $u=u(\varepsilon,y)\in D $, для которого


$\langle \nabla f(y),u\rangle \geq


\sup\limits_{x\in D}\langle \nabla f(y),x\rangle -\varepsilon$.


\item[(HU):] $\forall \varepsilon >0$ $\forall u \in D$ $\forall \zeta \in R:


\inf(f,D)\leq \zeta \leq \sup(f,D)$


можно приближенно решить задачу уровня (\ref{eq:m13}) при $\zeta=f(z)$,


т.\,е. отыскать элемент $\omega \in R^{n}: f(\omega)=\zeta $ такой, что


$\langle \nabla f(\omega ),u-\omega \rangle + \varepsilon \geq


\sup\limits_{v}\{ \langle \nabla f(v),u-v \rangle :\,\, f(v)=\zeta \}$.


\end{itemize}


Отметим, что для квадратичной функции задача уровня решается


аналитически~\cite{l7}.





Нетрудно видеть, что последовательное решение задач (\ref{eq:m12}) и 


(\ref{eq:m13}) не приводит, вообще говоря, к решению задачи (\ref{eq:m11})


совместной максимизации по $x$ и $y$. Поэтому возникает третья задача ---


задача выбора семейства точек $y^{i}$, $i=1,\dotsc,N$, на 


поверхности уровня $f(x)=f(z)\triangleq\zeta $, или, как будет


говориться в дальнейшем, задача аппроксимации поверхности уровня.





Пусть даны некоторый элемент $z\in D$ и число


$\varepsilon >0$. Обозначим $\zeta\triangleq f(z)$.


Рассмотрим аппроксимацию


\begin{equation}


\Re (\zeta)=\{ v^{1},\dotsc,v^{N} \in R^{n} : \, \,


f(v^{i})=\zeta,\ i=1,\dotsc N, \ N=N(\zeta) \}. 


\label{eq:m14}


\end{equation}


Согласно предположению (HL) $ \forall i=1,\dotsc,N$ определим $u^{i} \in D$


так, чтобы


\begin{equation}


\langle \nabla f(v^{i}),u^{i}\rangle \geq


\sup_{x\in D} \langle \nabla f(v^{i}),x \rangle - \varepsilon,


\label{eq:m15}


\end{equation}


а затем согласно предположению (HU) $\forall i=1,\dotsc,N$ определим 


элементы $\omega^{i} \in R^{n} : f(\omega^{i})=f(z)$ из условия


\begin{equation}


\langle \nabla f(\omega^{i}),u^{i}-\omega^{i} \rangle + \varepsilon \geq


\sup_{v}\{ \langle \nabla f(v),u^{i}-v \rangle :\,\, f(v)=\zeta \}.


\label{eq:m16}


\end{equation}





\begin{defn}


Набор $\Re(\zeta)$ назовем сильно $\varepsilon$-разрешающим,


если из того, 


что $z$ не является $\varepsilon$-решением задачи (\ref{eq:m1}), т.\,е.


\begin{equation}


f(z)< \sup(f,D)-\varepsilon,


\label{eq:m17}


\end{equation}


следует выполнение неравенств


\begin{align}


\text{i)}\ \; \eta(\zeta,\varepsilon )&\triangleq


\langle \nabla f(\omega^{j}),u^{j}-\omega^{j} \rangle =


\max_{1\leq i\leq N}\langle \nabla f(\omega^{i}),u^{i}-\omega^{i} \rangle >0;


\label{eq:m18} \\


\text{ii)}\ \; \eta(\zeta,\varepsilon )&\geq 


\langle \nabla f(y),p-y \rangle - 2\varepsilon,


\label{eq:m19}


\end{align}


где $p$ --- некоторый допустимый элемент $(p \in D)$, для которого


\begin{equation}


f(p)\geq f(z) + \theta [\sup(f,D)-f(z)]-\varepsilon,


\label{eq:m20}


\end{equation}


а число $\theta$ не зависит от $z$ и $\varepsilon$. При этом


$0<\theta <1$, и элемент $y\in R^{n}$ определяется из условий $f(y)=f(z)$,


\begin{equation}


\langle\nabla f(y),p-y\rangle \ge \|\nabla f(y)\|\, \|p-y\|- \varepsilon.


\label{eq:m21}


\end{equation}


\end{defn}





\begin{note}


Если найдется $p$, удовлетворяющее 


(\ref{eq:m20}), то существование $y:\, f(y)=f(z)$, удовлетворяющего


(\ref{eq:m21}), вытекает из предположения (HU) о решении


задачи (\ref{eq:m13}), т.\,к.


$$


\langle \nabla f(v),p-v \rangle \leq 


\| \nabla f(v) \|\, \| p-v \| \ \; \forall v \in Q.


$$


\end{note}





\begin{lem}\label{Lemma1} Пусть существует допустимый элемент


$p\in D: \; f(p)\leq f(z)+2\varepsilon $, для которого выполнено


$(\ref{eq:m20})$. Тогда из неравенства


$(\ref{eq:m18})$ следует $(\ref{eq:m19})$.


\end{lem}





{\bf Доказательство.} $\forall v: \; f(v)=f(z)$ имеем 


$$


\langle \nabla f(v),p-v \rangle \leq f(p)-f(v)=f(p)-f(z)\leq 2\varepsilon.


$$


Кроме того, из (\ref{eq:m18}) следует


$$


\langle \nabla f(v),p-v \rangle - 2\varepsilon \leq 0 <


\eta(\zeta,\varepsilon),


$$


т.\,е. (\ref{eq:m19}) тоже имеет место, что и требовалось.





\begin{defn}


Пусть заданы числа $\varepsilon >0$ и $\zeta$. Рассмотрим аппроксимацию


$$


W (\zeta) =\{ v^{1},\dotsc,v^{N} \in R^{n} :


f(v^{i})=\zeta,\ i=1,\dotsc, N,\ N=N(\zeta) \}.


$$


Построим точки $u^{i} \in D $ и $w^{i}: \, f(w^{i})=\zeta$ в соответствии


с формулами (\ref{eq:m15}) и (\ref{eq:m16}). Будем называть набор


$W (\zeta)$ слабо (или просто) $\varepsilon$-разрешающим, 


если из неравенства (\ref{eq:m17}) следует только (\ref{eq:m18}).


\end{defn}





\begin{prop}\label{Prop2} Если $\forall \varepsilon >0$ и


$\forall \lambda \in R$, \; $\inf(f,D) \leq \lambda < \sup (f,D)$,


существует процедура построения слабо $\varepsilon$-разрешающего набора


$W (\lambda, \varepsilon)$, то


$$


\forall z\in D, \ \ \sup(f,D) - \varepsilon > f(z) \geq \inf (f,D),


$$ 


можно построить и сильно $\varepsilon$-разрешающий набор


$\Re (\zeta)$, где $\zeta=f(z)$.


\end{prop}





{\bf Доказательство.} Пусть $ z\in D$, $f(z) \geq \inf (f,D)$,


$\varepsilon >0$.


Предположим, что $z$ не является $\varepsilon$-решением


задачи (\ref{eq:m1}), т.\,е. справедливо (\ref{eq:m17}). Положим


\begin{equation}


\lambda\triangleq f(z)+\theta[\sup(f,D)-f(z)],


\label{eq:m22}


\end{equation}


где число $\theta$ таково, что $0< \theta <1$. Очевидно,


$$


\inf(f,D) \leq f(z) < \lambda < \sup (f,D).


$$


Выберем число $\delta > 0$ по формуле


\begin{equation}


\delta \triangleq\frac{1- \theta}{2}[\sup(f,D)-f(z)].


\label{eq:m23}


\end{equation}


Нетрудно видеть, что в силу (\ref{eq:m22}) и (\ref{eq:m23})


на уровне $\lambda$ не существует $\delta$-решений задачи (\ref{eq:m1}),


т.\,е. $\lambda < \sup(f,D)- \delta$.


Поэтому по предположению можно построить слабо $\delta$-разрешающий набор 


$ W_{1}=W(\lambda)$.


Тогда из определения 2 следует, что найдутся $u^{j} \in D$, 


$w^{j}:\; f(w^{j})=\lambda$, при которых


$$


\eta (\lambda,\delta)\triangleq


\langle \nabla f(w^{j}),u^{j}-w^{j} \rangle >0.


$$ 


Отсюда с учетом выпуклости $f(\cdot)$ следует


\begin{equation}


f(u^{j})>f(w^{j})=\lambda\triangleq f(z)+\theta[\sup(f,D)-f(z)]>


f(z)+\theta[\sup(f,D)-f(z)]-\varepsilon=\lambda-\varepsilon.


\label{eq:m24}


\end{equation}


В частности, $f(u^{j})>f(z)$. Положим $p\triangleq u^{j} \in D$.


И вновь по предположению можно построить слабо $\varepsilon$-разрешающий


набор $W_{2}=W(\zeta)$. Рассмотрим множество


$$


\Re_{0}=W_{2}\cup \{ y \},


$$


где $y$ таково, что $f(y)=f(z)$ и при данных $p$ и $y$ cправедливо


(\ref{eq:m21}). Тогда из (\ref{eq:m24}) сразу же следует справедливость для


$p=u^j$ неравенства (\ref{eq:m20}) из определения 1.


С другой стороны, из предположения (HL) следует существование элемента


$u \in D$ такого, что 


$$


\langle \nabla f(y),u \rangle \geq 


\sup_{x \in D}\langle \nabla f(y),x \rangle - \varepsilon.


$$


Отсюда вытекает


\begin{equation}


\langle \nabla f(y),u-y \rangle \geq


\langle \nabla f(y),p-y \rangle - \varepsilon.


\label{eq:m25}


\end{equation}





Далее, по предположению (HU) можно построить $w \in R^{n}: \, f(w)=f(z)$,


для которого


$$


\langle \nabla f(w),u-w \rangle + \varepsilon \geq


\sup_{v} \{ \langle \nabla f(v),u-v \rangle : \, f(v)=f(z) \} \geq


\langle \nabla f(y),u-y \rangle.


$$


Отсюда с помощью (\ref{eq:m25}) имеем


$$


\langle \nabla f(w),u-w \rangle \geq


\langle \nabla f(y),p-y\rangle -2\varepsilon.


$$ 


А тогда для аппроксимации $\Re_{0}=W_{2} \cup \{ y \}$ получаем


$$


\eta(\zeta,\varepsilon) \geq \langle \nabla f(w),u-w \rangle \geq


\langle \nabla f(y),p-y\rangle -2\varepsilon .


$$


Итак, неравенство (\ref{eq:m19}) из определения 1 оказывается выполненным.


А неравенство (\ref{eq:m18}) справедливо по построению


$W_{2}=W(\zeta)$. Поэтому $\Re_{0}$ является сильно


разрешающим набором.





\section{$\Re$-стратегия и ее сходимость}





Пусть заданы элемент $x_{0} \in D$ и числовая последовательность 


$\{ \varepsilon_{k} \}: \,


\varepsilon_{k}>0$, $k=0,1,2,\dotsc$,


$\varepsilon_{k} \downarrow 0$ ($k \rightarrow \infty$).


Предположим, что найдется константа $\varkappa> 0$, для которой 


\begin{equation}


\| \nabla f(v) \| \geq \varkappa >0 \ \; \forall v: \,


f(v) \geq f(x^{0}). 


\label{eq:m26}


\end{equation}





Кроме того, рассмотрим предположение


\begin{itemize}


\item[(HR):] $\forall \varepsilon >0$ и для любой $\varepsilon$-стационарной


точки $z \in D$,


не являющейся $\varepsilon$-решением задачи (\ref{eq:m1}),


можно построить слабо разрешающий набор $W(f(z))$.


\end{itemize}


Согласно предложению 2 предположение (HR) гарантирует возможность построения


сильно разрешающего набора для любой допустимой точки, не являющейся


$\varepsilon$-глобальным решением. Легко видеть, что для любой нестационарной 


допустимой точки $x$ множество $W(f(x))=\{ x \}$ является слабо разрешающим


набором благодаря выпуклости $f(\cdot)$. Действительно, в этом случае 


найдется $u=u(x) \in D $, для которого


$$


0< \langle \nabla f(x),u-x \rangle \leq 


\sup_{v} \langle \nabla f(v),u-v \rangle ,


$$ 


что, как нетрудно видеть, согласуется с определением 2.





Для решения задачи (\ref{eq:m1}) предлагается следующий алгоритм.


\begin{itemize}


\begin{itemize}


\item[Шаг 0.] Полагаем $k:=0$. Пусть задана некоторая допустимая


точка $x^{0}$.


\item[Шаг 1.] Исходя из точки $x^{k}$, методом локального подъема пункта~2


получаем $\varepsilon_{k}$-ста\-ци\-о\-нар\-ную точку $z^{k}$,


$f(z^{k})\triangleq\zeta_{k}$, т.\,е. точку, удовлетворяющую


условию


$$


\sup_{x\in D} \langle \nabla f(z^k),x-z^k \rangle \leq \varepsilon_k .


$$ 


\item[Шаг 2.] Строим сильно $\varepsilon_k$-разрешающий набор


$$


\Re_{k}=\Re(\zeta_k)=\{ v^{1},\dotsc ,v^{N}\mid


f(v^{i})=\zeta_{k},\ i=1,\dotsc,N,\ N=N(k) \}.


$$


\item[Шаг 3.] Для $ i=1,\dotsc,N$ находим точки $u^{i} \in D $ по правилу


(\ref{eq:m15}).


\item[Шаг 4.] Для $i=1,\dotsc ,N$ определяем точки 


$w^{i}: \, f(w^{i})=\zeta_{k} $ по правилу (\ref{eq:m16}), 


где $\zeta = \zeta_{k}$.


\item[Шаг 5.] Формируем величину


$\eta_{k}=\eta(\zeta_{k},\varepsilon_{k})=


\langle \nabla f(w^{j}),u^{j}-w^{j} \rangle \triangleq


\max\limits_{i} \langle \nabla f(w^{i}),u^{i}-w^{i} \rangle$. 


\item[Шаг 6.] Если $\eta_{k}>0$, то полагаем $x^{k+1}:=u^{j}$, $k:=k+1$


и возвращаемся на шаг~1.


\item[Шаг 7.] Если $\eta_{k} \leq 0$, то полагаем $x^{k+1}:=z^{k}$, $k:=k+1$


и возвращаемся на шаг~1.


\end{itemize}


\end{itemize}





\begin{note}


Шаги 1, 2, 3, 4 $\Re$-алгоритма являются обоснованными


ввиду предложения 1 и предположений (HR) (HL) и (HU) соответственно.


\end{note}


 


\begin{note}


Из описания шага~1 видно, что последовательность,


генерируемая $\Re$-ал\-го\-рит\-мом, является последовательностью 


$\varepsilon_{k}$-стационарных точек.


\end{note}





\begin{note}


При практической реализации алгоритма останов производится, если


$\eta_{k}\leq 0 $ и $\varepsilon_{k} \leq \delta $, где $\delta$ 


--- заданная точность. В этом случае из определения сильно разрешающего


набора следует


$$


f(z^{k}) \geq \sup(f,D) - \varepsilon_{k} \geq \sup(f,D)- \delta,


$$


так что $z^{k}$ является $\delta$-решением.


\end{note}





\begin{note}


Нетрудно видеть, что вышеописанный метод не является


алгоритмом в общепринятом смысле, поскольку, например, не уточнено, 


каким алгоритмом решать линеаризованную задачу или задачу уровня. Не


указано также, как строить (сильно или слабо) разрешающий набор. Да


и на первом шаге пользователь может выбрать другой метод локального 


подъема, если он считает его более подходящим в том или ином случае. 


Отметим только, что эти методы должны работать очень быстро, поскольку


они применяются по несколько раз на каждой итерации.


\end{note}





Поэтому разумнo назвать представленный выше метод глобального


поиска принципиальным алгоритмом \cite{l23}


или стратегией глобального поиска \cite{l9}. 


Поскольку понятие разрешающего набора является определяющим для


успешного поиска глобального максимума в задаче (\ref{eq:m1}), то будем


называть вышеописанный метод для краткости $\Re$-стратегией.





\begin{thm}\label{Theorem2}


Пусть выполнены предположения $(\ref{eq:m2})$, $(\ref{eq:m26})$,


$\mathrm{(HR)}$, $\mathrm{(HL)}$ и 


$\mathrm{(HU)}$. Тогда последовательность $\{ z^{k} \}$, генерируемая


$\Re$-стратегией, является максимизирующей в задаче $(\ref{eq:m1})$.


\end{thm}





{\bf Доказательство.}


а) Пусть $\eta_{k} \leq 0$ $\forall k \geq m \geq 0$. 


Тогда из определения разрешающего набора


$f(z^{k}) \geq \sup(f,D)-\varepsilon_{k}$.


Поэтому $\lim\limits_{k \rightarrow \infty}f(z^{k})=\sup(f,D)$.





б) Пусть теперь $\eta_{k}>0$ $\forall k=0,1,2,\dotsc$ По построению


$z^{k+1}$ является $\varepsilon_{k+1}$-стационарной точкой,


полученной локальным подъемом раздела 2, исходя из


$x^{k+1}:=u^{j_{k}} $, где


$$


\eta_{k}\triangleq \eta(f(z^{k}),\varepsilon_{k})=


\langle \nabla f(w^{j_{k}}),u^{j_{k}}-w^{j_{k}} \rangle = 


\max_{1\leq i \leq N} \langle \nabla f(w^{i}),u^{i}-w^{i} \rangle >0.


$$ 


Поэтому в силу выпуклости $f(\cdot)$ имеем


\begin{equation}


0< \eta_{k} \leq f(u^{j_{k}})-f(w^{j_{k}}) \leq f(z^{k+1})-f(z^{k}).


\label{eq:m27}


\end{equation}


Нетрудно видеть, что при этом


$f(z^{k})<f(x^{k+1}) \leq f(z^{k+1})$.





Таким образом, числовая последовательность


$\{ f(z^{k}) \}$, $k=0,1,2,\dotsc$, является монотонно возрастающей.


И поскольку согласно (\ref{eq:m2}) $f(\cdot)$ ограничена сверху


на $D$, то существует конечный предел


$\lim\limits_{k \to \infty}f(z^{k})=\ov{f}$.


При этом из (\ref{eq:m27}) следует, что


\begin{equation}


\eta_{k} \downarrow 0 \ \;(k\to +\infty).


\label{eq:m28}


\end{equation}


Предположим, тем не менее, что $\{ z^{k} \}$ не является максимизирующей,


т.\,е. найдется число $\gamma >0$, при котором 


\begin{equation}


\zeta_{k}\triangleq f(z^{k})\leq \ov{f} < \sup (f,D)-\gamma.


\label{eq:m29}


\end{equation}





Поскольку в каждой точке $z^{k}$ можно построить слабо разрешающий набор,


а потому согласно предложению 2 и сильно разрешающий набор,


то будем иметь 


\begin{equation}


\eta_{k}=\eta(z^{k},\varepsilon_{k})\geq 


\langle \nabla f(y^{k}),p^{k}-y^{k} \rangle -2\varepsilon_{k}, 


\label{eq:m30}


\end{equation}


где $p^{k}$ таково, что 


$f(p^{k})\geq f(z^{k})+\theta[\sup(f,D)-\zeta_{k}]-\varepsilon_{k}$.


Из последнего неравенства и (\ref{eq:m29}) вытекает


\begin{equation}


f(p^{k})\geq f(z^{k})+\theta \gamma -\varepsilon_{k},


\label{eq:m31}


\end{equation}


откуда в свою очередь следует


\begin{equation}


\liminf_{k \to \infty}f(p^{k})\geq \ov{f}+\theta \gamma.


\label{eq:m32}


\end{equation}





С другой стороны, поскольку $\theta \gamma -\varepsilon_{k} >0 $ при


достаточно больших $k$, то из (\ref{eq:m31}) выводим 


$f(p^{k})>f(z^{k})\triangleq\zeta_{k}$ $\forall k\geq k_{0}$.


Выберем теперь $y_{k},\:f(y^{k})=\zeta_{k}$, таким образом, чтобы 


$\langle \nabla f(y^{k}),p^{k}-y^{k} \rangle \geq


\| \nabla f(y^{k}) \|\, \|p^{k}-y^{k} \| -\varepsilon_{k}$.


С учетом последнего неравенства из (\ref{eq:m30}) следует


$\eta_{k}+2\varepsilon_{k} \geq 


\| \nabla f(y^{k}) \|\, \|p^{k}-y^{k} \| -\varepsilon_{k}$.


Отсюда получаем в силу предположения (\ref{eq:m26})


$\eta_{k}+3\varepsilon_{k} \geq \varkappa \|p^{k}-y^{k} \|$.


Поэтому с помощью (\ref{eq:m28}) заключаем, что 


$\|p^{k}-y^{k} \| \downarrow 0$ ($k \to \infty$).


А тогда в силу непрерывности $f(\cdot)$ имеем


$$


\lim_{k \to \infty}f(p^{k})=\lim_{k \to \infty}f(y^{k})=


\lim_{k \to \infty}\zeta_{k}=\ov{f}.


$$


Однако последние равенства несовместны с (\ref{eq:m32}), поскольку 


$\theta >0, \gamma >0 $.


Итак, предположение о том, что последовательность $\{ z^{k} \}$ не является


максимизирующей, привело нас к абсурду. Следовательно, $\{ z^{k} \}$ 


оказывается максимизирующей.





в) Очевидно, в общем случае последовательность $\{ \eta_{k} \}$


разбивается на две


подпоследовательности $\{ \eta_{k_{s}} \}$ и $ \{ \eta_{k_{t}} \}$


посредством условий


$\eta_{k_{s}}\leq 0 $ и $\eta_{k_{t}} >0$ 


так, что ${k_{s}}\cup {k_{t}}=\{0, 1, 2, \dotsc \}.$


Cоответственно и последовательность $ \{ z^{k} \}$ распадается на две


подпоследовательности $\{z^{k_s}\}$ и $\{z^{k_t}\}$.


%% $\{ z^{k} \}=\{z^{k_s}\} \cup \{z^{k_t}\}$.


При этом $\{z^{k_s}\}$, для которой $\eta_{k_{s}}\leq 0$, $s=0,1,2,\dotsc$,


будет максимизирующей по части~a), а подпоследовательность


$\{z^{k_t}\}$ будет максимизирующей по части~б) доказательства. 


Поэтому и вся последовательность $\{ z^{k} \}$


будет максимизирующей в задаче (\ref{eq:m1}).





\begin{thebibliography}{99}


\bibitem{l1}


Васильев Ф.П. {\em Численные методы решения экстремальных


задач.} -- М.: Наука, 1988. -- 552~с.


\bibitem{l23}


Полак Э. {\em Численные методы оптимизации. Единый подход.}


-- М.:Мир, 1974. -- 374 с.


\bibitem{l2}


Поляк Б.Т. {\em Введение в оптимизацию.} -- М.: Наука,


1983. -- 384 с.


\bibitem{l3}


Туй Х. {\em Вогнутое программирование при линейных


ограничениях} // ДАН СССР. -- 1964. -- Т.\,159. -- \No\,9. -- С.\,32--35.


\bibitem{l4}


Horst R., Tuy H. {\em Global optimization $($deterministic 


approaches$)$.} -- 2-nd edition. -- Berlin: Springer-Verlag, 1993. -- 698 p.


\bibitem{l5}


Horst R., Pardalos P.M., Thoai N.V. {\em Introduction to


global optimization.} -- Dordrecht: Kluwer Acad. Publ.,


1995. -- 246 p.


\bibitem{l8}


Hiriart-Urruty J.-B., Ledyaev Y.A. {\em Note on the


characterization


of the global maxima of a $($tangentially$)$ convex function over a 


convex set} // J. of Convex Anal. -- 1996. -- V.\,3. -- \No\,1.


-- P.\,55--61.


\bibitem{l6}


Стрекаловский А.С. {\em К проблеме глобального экстремума в 


невыпуклых задачах оптимизации} // Изв. вузов. Математика. -- 1990. -- \No\,8.


-- С.\,74--80.


\bibitem{l7}


Стрекаловский А.С. {\em О поиске глобального максимума выпуклого


функционала на допустимом множестве} // Журн. вычисл. матем. и матем.


физ. -- 1993. -- Т.\,33. -- \No\,3. -- С.\,349--364.


\bibitem{l9}


Strekalovsky A.S., Tsevendorj I. {\em Testing the $R$-strategy


for a reverse convex problem} // J. of Global Optimization. -- 1998. --


\No\,13. -- P.\,61--74.


\bibitem{l10}


Кузнецова А.А., Стрекаловский А.С., Цэвээндорж И. 


{\em Об одном подходе к решению целочисленных задач оптимизации} //


Журн. вычисл. матем. и матем. физ. -- 1999. -- Т.\,39. -- \No\,1. --


С.\,9--16. 





\end{thebibliography}





\vskip 2cm





\postup{\em Институт динамики систем и теории}{$\qquad$\it Поступили}


\postup{\em управления Сибирского отделения}{{\it первый вариант} 21.04.1998}


\postup{\em Российской Академии Наук}{{\it окончательный вариант} 15.06.1999}





\label{end}


\end{document}


